\chapter{Noise, Sensitivity, and Dynamic Range}
\label{ch:noise}
A receiver is judged at two extremes: the weakest signal it can hear and the strongest
it can tolerate. Both are numbers this book can now derive rather than quote, because
both come from material already in hand---the noise floor from thermal physics plus
\cref{ch:resistive}'s maximum-power-transfer condition, and the strong-signal limits
from the same nonlinearity \cref{sec:mixers} used to explain mixing.

That is the shape of the chapter. The lower limit is set by \emph{noise}: how much of
it the receiver adds, and where in the chain that addition happens. The upper limit is
set by \emph{nonlinearity}: how gracefully the front end fails when overdriven. The gap
between them is the \textbf{dynamic range}, and the several numbers that quantify it
turn out to be consequences of a single power series.

\section{The Noise Floor: Where \texorpdfstring{\(-174\)}{-174} dBm/Hz Comes From}
\label{sec:noisefloor}
Every resistor at a temperature above absolute zero generates a small random voltage,
because its charge carriers are in thermal motion. Johnson and Nyquist showed the
open-circuit noise voltage across a resistance \(R\) in a bandwidth \(B\) is
\[
v_{\RMS}=\sqrt{4kTRB},
\]
where \(k=\SI{1.38e-23}{\joule\per\kelvin}\) is Boltzmann's constant and \(T\) is the
absolute temperature. That is the one physical input this chapter needs; everything
below is circuit theory.

What a receiver actually cares about is not the voltage but the \emph{power} that
resistor can deliver to it. And \cref{ch:resistive} already settled how much that
is---maximum power transfer occurs into a matched load \(R\), which by the voltage
divider takes half the source voltage. So the \textbf{available noise power} is
\[
P_n=\frac{\bigl(v_{\RMS}/2\bigr)^2}{R}
=\frac{4kTRB}{4R}
\qquad\Longrightarrow\qquad
\boxed{\;P_n=kTB\;}
\]
Notice that the resistance cancels completely: available noise power depends only on
temperature and bandwidth, not on the resistor generating it. And notice the factor of
\(4\) in the Johnson formula---it is there precisely so that this result comes out
clean, which is the usual sign that the available power, not the open-circuit voltage,
is the physically meaningful quantity.

Now evaluate it. At a reference temperature \(T_0=\SI{290}{\kelvin}\),
\[
kT_0=(1.38\times10^{-23})(290)=\SI{4.00e-21}{\watt\per\hertz},
\]
and converting to dBm (\cref{sec:sunits}) by dividing by \SI{1}{\milli\watt},
\[
10\log_{10}\frac{4.00\times10^{-21}}{10^{-3}}
=10\log_{10}(4.00\times10^{-18})
\approx\boxed{\;\SI{-174}{dBm}\ \text{per hertz}\;}
\]
That is the number every receiver specification is measured against. It is not a
property of any circuit---it is what a matched resistor at room temperature hands you,
and therefore the best any receiver could possibly do.

The bandwidth dependence needs no further work: since \(P_n\propto B\), widening the
bandwidth raises the floor by \(10\log_{10}B\). So
\[
\text{floor (dBm)}=-174+10\log_{10}B .
\]
A \SI{500}{\hertz} CW filter sits \(10\log_{10}500\approx\SI{27}{\decibel}\) above the
per-hertz figure, at \SI{-147}{dBm}; a \SI{2.4}{\kilo\hertz} SSB filter sits
\SI{34}{\decibel} up, at \SI{-140}{dBm}. \emph{Narrower is quieter}, and by an amount
you can compute.

\begin{workedbox}[: from microvolts to dBm]
Receiver sensitivity is quoted both ways, so the bridge is worth having. In a
\SI{50}{\ohm} system, \(P=V_{\RMS}^2/R\), so
\[
P_{\mathrm{dBm}}=10\log_{10}\frac{V^2/50}{10^{-3}}
=20\log_{10}V+13.0 ,
\]
with \(V\) in volts---or, with \(V\) in the microvolts receivers are usually specified
in, \(P_{\mathrm{dBm}}=20\log_{10}V_{\mu\mathrm{V}}-107\). Check it on a common
figure: \(\SI{0.5}{\micro\volt}\) gives
\(20\log_{10}(0.5)-107=-6-107=\SI{-113}{dBm}\). So ``\SI{0.5}{\micro\volt}
sensitivity'' and ``\SI{-113}{dBm}'' are the same claim.
\end{workedbox}

\begin{unitsbox}
\[
kTB=\left(\frac{\si{\joule}}{\si{\kelvin}}\right)(\si{\kelvin})
\left(\frac{1}{\si{\second}}\right)=\frac{\si{\joule}}{\si{\second}}=\si{\watt}.
\]
\end{unitsbox}

\section{Noise Figure and Noise Temperature}
\label{sec:noisefigure}
A real receiver adds noise of its own. The standard measure is the \textbf{noise
figure}: the factor by which the receiver degrades the signal-to-noise ratio,
\[
F=\frac{\mathrm{SNR}_{\mathrm{in}}}{\mathrm{SNR}_{\mathrm{out}}},
\]
with the input driven by a source at \(T_0\). An ideal, noiseless amplifier multiplies
signal and noise equally and leaves the ratio alone, so \(F=1\) (\SI{0}{\decibel}); any
real one has \(F>1\).

It is often more convenient to express the same thing as an \textbf{equivalent noise
temperature} \(T_e\)---the temperature a resistor at the input would need in order to
account, all by itself, for the noise the device actually adds. Writing the total output
noise of a stage of gain \(G\) as
\[
N_{\mathrm{out}}=G\,kB\,(T_0+T_e)
\]
and comparing with the ideal \(GkBT_0\) gives
\[
\boxed{\;F=1+\frac{T_e}{T_0}\;}
\]
so the two descriptions carry identical content.

This also explains a step usually presented as a rule. Since \(F\) multiplies a power
ratio, and decibels turn multiplication into addition (\cref{sec:cascadeadd}), the
receiver's actual noise floor is the thermal floor \emph{plus} the noise figure:
\[
\text{floor}=-174+10\log_{10}B+\mathrm{NF}_{\dB}.
\]
Nothing is being asserted---the addition is what taking a logarithm of a product does.

\section{Cascades: Why the First Stage Decides}
\label{sec:friis}
Receivers are chains, so the useful question is how the stages combine. Take two, with
gains \(G_1,G_2\) and excess noise temperatures \(T_{e1},T_{e2}\). Stage 1 delivers
\(G_1kB(T_0+T_{e1})\) to stage 2, which amplifies all of it and adds its own:
\[
N_{\mathrm{out}}=G_1G_2\,kB\!\left(T_0+T_{e1}+\frac{T_{e2}}{G_1}\right).
\]
The structure is the whole result: stage 2's contribution is \emph{divided by the gain
ahead of it}. Converting to noise figures gives \textbf{Friis's formula},
\[
\boxed{\;
F=F_1+\frac{F_2-1}{G_1}+\frac{F_3-1}{G_1G_2}+\cdots\;}
\]
Three consequences follow, and none of them is obvious without the derivation:
\begin{description}
\item[The first stage dominates.] Every later stage is discounted by the accumulated
gain in front of it. A receiver's noise figure is very nearly its first amplifier's,
provided that amplifier has useful gain.
\item[A loss ahead of the receiver costs its full value.] A passive attenuator or a
length of feed line with loss \(L\) has \(F=L\) and \(G=1/L\), so \(n\) decibels of
line loss adds \(n\) decibels directly to the system noise figure. There is no gain
ahead of it to discount it.
\item[Therefore put the preamplifier at the antenna.] At the mast, the preamp is stage
one: its gain then discounts both the feed line's loss and the receiver's own noise. In
the shack, the line's loss comes first and has already been added in full, and no amount
of subsequent gain removes it.
\end{description}

\begin{controlsbox}
\textbf{Two chains, two different rules.} \Cref{sec:sunits}'s link budget adds gains and
subtracts losses in decibels, treating every stage alike. This chain \emph{weights}
them, discounting each stage by the gain preceding it. The difference is that a link
budget tracks the \emph{signal}, which every stage multiplies equally, while Friis tracks
the \emph{ratio} of signal to a noise that each stage adds afresh. Same cascade, two
different questions.
\end{controlsbox}

\section{Sensitivity and MDS}
The two results above define the weakest usable signal. The \textbf{minimum discernible
signal} (MDS) is conventionally the input level that arrives at the same strength as the
noise---a signal-to-noise ratio of \SI{0}{\decibel}:
\[
\mathrm{MDS}=-174+10\log_{10}B+\mathrm{NF}_{\dB}.
\]
It is measured on the bench by watching for a \SI{3}{\decibel} rise in output as the
signal is applied, and that number is not arbitrary either: when signal power equals
noise power, the total is twice the noise alone, and \(10\log_{10}2\approx\SI{3}{\decibel}\).

\begin{workedbox}[: the sensitivity of a real receiver]
A receiver has \(\mathrm{NF}=\SI{8}{\decibel}\) and a \SI{500}{\hertz} CW filter. Then
\[
\mathrm{MDS}=-174+10\log_{10}500+8
=-174+27+8=\SI{-139}{dBm},
\]
which by the bridge above is about \SI{0.025}{\micro\volt}. Switch to a
\SI{2.4}{\kilo\hertz} SSB filter and the floor rises to \(-174+34+8=\SI{-132}{dBm}\)
--- \SI{7}{\decibel} worse, purely from bandwidth. And improving the noise figure from
\SI{8}{\decibel} to \SI{3}{\decibel} would buy \SI{5}{\decibel}: on HF, where external
atmospheric noise usually dominates anyway, the filter choice matters more than the
preamplifier.
\end{workedbox}

\section{One Power Series Generates Every Strong-Signal Number}
\label{sec:powerseries}
Now the other end. \Cref{sec:mixers} said that real mixers and overdriven amplifiers ``do
not multiply cleanly'' and left it there. Making that precise turns out to produce every
strong-signal specification at once.

Model a weakly nonlinear stage by expanding its transfer characteristic as a power
series about the operating point:
\[
y=a_1x+a_2x^2+a_3x^3+\cdots
\]
The first term is the wanted linear gain; the rest are the trouble. Drive it with two
equal tones, \(x=A(\cos\omega_1t+\cos\omega_2t)\), and expand using the same
product-to-sum identity as \cref{sec:mixers}. Three results fall out of one calculation.

\subsection{The Slopes Are the Exponents}
The \(a_2x^2\) term produces components at \(\omega_1\pm\omega_2\) with amplitude
proportional to \(a_2A^2\); the \(a_3x^3\) term produces the third-order products
\(2\omega_1-\omega_2\) and \(2\omega_2-\omega_1\) with amplitude
\(\tfrac34a_3A^3\). So if the input amplitude rises by a factor \(r\):
\begin{itemize}
\item the wanted output rises as \(r\) --- \(\SI{1}{\decibel}\) per \si{\decibel};
\item second-order products rise as \(r^2\) --- \(\SI{2}{\decibel}\) per \si{\decibel};
\item third-order products rise as \(r^3\) --- \(\SI{3}{\decibel}\) per \si{\decibel}.
\end{itemize}
\emph{The famous slopes are just the exponents of the series.} There is nothing to
memorize once the model is written down, and the reason third-order products are the
dangerous ones is already in \cref{sec:mixers}: they land next to the signals that
created them.

\subsection{Gain Compression}
\label{sec:compression}
The cubic term does something else too. Collecting the components at \(\omega_1\)
itself---the cube of a sum contributes back to the fundamental---the effective gain
becomes
\[
a_1+\tfrac94a_3A^2
\qquad\text{(two equal tones; }\tfrac34\text{ for a single tone).}
\]
In a real device \(a_3\) has the \emph{opposite} sign to \(a_1\), so the gain
\emph{falls} as drive increases. That single term is gain compression, blocking, and
desensitization---three names for one coefficient. The \textbf{\(\SI{1}{\decibel}\)
compression point} is where the gain has dropped by \SI{1}{\decibel}, i.e.\ where the
bracket equals \(10^{-1/20}a_1\approx0.891\,a_1\).

\subsection{The Third-Order Intercept}
Finally, compare the third-order product with the wanted output. The product grows three
times as fast, so extrapolating both straight lines on a log--log plot they must cross.
Setting \(\tfrac34|a_3|A^3=a_1A\) gives the crossing amplitude,
\[
A_{\mathrm{IIP3}}^2=\frac{4}{3}\left|\frac{a_1}{a_3}\right|,
\]
and the input power there is the \textbf{input third-order intercept} \(\mathrm{IIP3}\).
Two things about it are worth saying plainly. First, it is a \emph{fiction}: the
intercept is an extrapolation to a power level that would destroy the receiver long
before reaching it, so IIP3 is best read as a linearity coefficient wearing the units of
power. Second, since both \(P_{\SI{1}{\decibel}}\) and IIP3 are set by the same ratio
\(|a_1/a_3|\), they are not independent---dividing the two expressions gives
\[
\mathrm{IIP3}\approx P_{\SI{1}{\decibel}}+\SI{9.6}{\decibel}
\]
for a purely cubic nonlinearity. A useful sanity check on a datasheet: if the two
numbers differ by wildly more or less than \SI{10}{\decibel}, something beyond the
cubic term is at work.

\section{Dynamic Range: Two Numbers, Two Different Failures}
\label{sec:dynamicrange}
The receiver is usable between the noise floor and whichever strong-signal effect bites
first. The two effects fail differently, so they get separate figures.

\textbf{Blocking dynamic range} uses compression: it is simply how far the
\SI{1}{\decibel} compression point sits above the MDS,
\[
\mathrm{BDR}=P_{\SI{1}{\decibel}}-\mathrm{MDS},
\]
a plain subtraction, \emph{because} compression is a one-for-one phenomenon.

\textbf{Intermodulation dynamic range} uses the third-order products, and here the
slope enters. From the slopes above, the third-order product referred to the input obeys
\[
P_{\mathrm{IM3}}=3P_{\mathrm{in}}-2\,\mathrm{IIP3}
\]
in decibels---which passes the obvious check: at \(P_{\mathrm{in}}=\mathrm{IIP3}\) it
returns \(\mathrm{IIP3}\), as the definition of an intercept requires. The usable limit
is the drive level at which those products just reach the noise floor, so set
\(P_{\mathrm{IM3}}=\mathrm{MDS}\) and solve:
\[
P_{\mathrm{in}}=\frac{2\,\mathrm{IIP3}+\mathrm{MDS}}{3}
\qquad\Longrightarrow\qquad
\boxed{\;\mathrm{DR3}=P_{\mathrm{in}}-\mathrm{MDS}
=\tfrac{2}{3}\bigl(\mathrm{IIP3}-\mathrm{MDS}\bigr)\;}
\]
\emph{The \(2/3\) is slope-3 geometry and nothing else.} The same argument at slope two
gives \(\mathrm{DR2}=\tfrac12(\mathrm{IIP2}-\mathrm{MDS})\)---the coefficient is always
\((n-1)/n\) for order \(n\), which is worth noticing because it makes the number
memorable instead of magic.

\begin{workedbox}[: a full receiver budget]
A receiver has \(\mathrm{NF}=\SI{10}{\decibel}\),
\(\mathrm{IIP3}=\SI{-10}{dBm}\), \(P_{\SI{1}{\decibel}}=\SI{-20}{dBm}\), measured in a
\SI{500}{\hertz} bandwidth. Then
\[
\mathrm{MDS}=-174+27+10=\SI{-137}{dBm},
\]
\[
\mathrm{BDR}=-20-(-137)=\SI{117}{\decibel},
\qquad
\mathrm{DR3}=\tfrac23\bigl(-10-(-137)\bigr)=\tfrac23(127)\approx\SI{85}{\decibel}.
\]
Two readings. The intermodulation figure is \SI{32}{\decibel} \emph{worse} than the
blocking figure, which is why DR3 is the number that characterizes a crowded band---it
is third-order products, not compression, that spoil a contest weekend. And the
\SI{10}{\decibel} gap between \(P_{\SI{1}{\decibel}}\) and IIP3 is close to the
\SI{9.6}{\decibel} the cubic model predicts, so these numbers are mutually consistent.
\end{workedbox}

\begin{workedbox}[: why the attenuator switch helps]
The front-panel attenuator looks like it can only make things worse, and the power
series shows why it does not. Insert \(L\) decibels of attenuation. The wanted signal
falls by \(L\); the third-order products, rising at \SI{3}{\decibel} per \si{\decibel},
fall by \(3L\). So the product-to-signal ratio \emph{improves} by
\[
3L-L=2L\ \si{\decibel} .
\]
Ten decibels of attenuation buys \SI{20}{\decibel} of intermodulation relief. The cost,
by \cref{sec:friis}, is \(L\) decibels of noise figure---which is free whenever external
noise already dominates the receiver's own, because then signal and noise are attenuated
together and the signal-to-noise ratio does not change at all. That is precisely the
situation on the low HF bands, and precisely why the switch exists.
\end{workedbox}

\section{Phase Noise and Reciprocal Mixing}
\label{sec:phasenoise}
One more mechanism raises the noise floor, and it is the most on-thesis of all of them:
it is purely a statement about where an oscillator's poles sit.

\Cref{ch:active} noted qualitatively that a high-\(Q\) resonator pins an oscillator's
pole pair tightly to the \(\jj\omega\) axis and gives a sharp spectral line. Make that
quantitative. Write the local oscillator with a small random phase disturbance
\(\phi(t)\) and expand for small \(\phi\):
\[
\cos\bigl(\omega_{\mathrm{LO}}t+\phi(t)\bigr)
\approx\cos\omega_{\mathrm{LO}}t-\phi(t)\sin\omega_{\mathrm{LO}}t .
\]
A pure carrier, \emph{plus} a term carrying the phase noise in quadrature with it. So a
jittering oscillator is not a line: it is a line with \textbf{skirts}, whose density at
an offset \(\Delta f\) from the carrier is quoted in dBc/Hz---decibels relative to the
carrier, per hertz.

Now put that oscillator in a mixer. The mixer multiplies, so a strong interfering signal
at offset \(\Delta f\) from the tuned frequency beats against the skirt component at
that same offset, and the product lands \emph{inside} the IF passband, where it is
indistinguishable from noise. Its level is
\[
P_{\text{noise}}=P_{\text{interferer}}+\mathcal{L}(\Delta f)+10\log_{10}B ,
\]
with \(\mathcal{L}\) the phase-noise density. This is \textbf{reciprocal mixing}, and
its signature is distinctive: the noise floor \emph{rises as you tune towards} a strong
signal, because the skirt is steeper closer in. Note that it is a noise-floor
mechanism, not a distortion mechanism---which is exactly why it is specified separately
from \cref{sec:dynamicrange}'s dynamic ranges.

\subsection{Where the Skirt's Shape Comes From}
The final step connects it to a pole. An oscillator's tank filters the active device's
own noise, and by \cref{ch:rlc} a resonator of loaded quality factor \(Q_L\) is a
band-pass whose half-bandwidth is \(\omega_0/2Q_L\). Above that corner the filtering
rolls off at \(-\SI{20}{\decibel}\) per decade, so
\[
\mathcal{L}(\Delta f)\ \propto\ \frac{1}{Q_L^2}\cdot\frac{1}{\Delta f^{2}} .
\]
Two readings, both practical. The \(-\SI{20}{\decibel}\)/decade slope is the
single-pole roll-off of \cref{ch:splane}, appearing here as the shape of a noise skirt.
And the \(1/Q_L^2\) means \textbf{doubling the loaded \(Q\) buys \SI{6}{\decibel} of
phase noise}---which finally puts a number on the tank-loading trade-off of
\cref{ch:crossproblems}: coupling the next stage loosely is not fussiness, it is
\SI{6}{\decibel} per octave of \(Q\).

\begin{controlsbox}
\textbf{Everything about an oscillator's spectrum is its pole location.} A pole exactly
on the \(\jj\omega\) axis is a perfect sinusoid---an infinitely thin spectral line. Real
loss and noise let the pole wander slightly, and the width of that wander \emph{is} the
skirt. High \(Q_L\) means a stiffly held pole, hence a narrow skirt, hence low phase
noise, hence less reciprocal mixing. One geometric fact explains an oscillator
specification, a receiver specification, and a design rule.
\end{controlsbox}

\section{Noise Bandwidth Is Not the \texorpdfstring{\(-3\)}{-3} dB Bandwidth}
\label{sec:noisebw}
One loose end in the arithmetic above: the \(10\log_{10}B\) term assumed a filter that
passes everything inside \(B\) and nothing outside. Real filters do not, so \(B\) needs
defining honestly.

Noise arrives at every frequency, so the noise a filter passes is the integral of its
power response. Define the \textbf{equivalent noise bandwidth} as the width of the ideal
rectangular filter that would pass the same noise:
\[
B_n=\frac{\displaystyle\int_0^\infty|H(f)|^2\,\dd f}{|H|^2_{\max}} .
\]
Evaluate it on the book's own first-order low-pass, \(|H|^2=1/[1+(f/f_c)^2]\) from
\cref{ch:rc}. The integral is a standard arctangent:
\[
\int_0^\infty\frac{\dd f}{1+(f/f_c)^2}=f_c\Bigl[\arctan\frac{f}{f_c}\Bigr]_0^\infty
=\frac{\pi}{2}f_c
\qquad\Longrightarrow\qquad
B_n=\frac{\pi}{2}f_c\approx1.57\,f_{\SI{3}{\decibel}} .
\]
So a single-pole filter is \(10\log_{10}1.57\approx\SI{1.96}{\decibel}\) noisier than its
nameplate bandwidth suggests---the skirts let in noise the \(-\SI{3}{\decibel}\) figure
does not account for. For a Butterworth response of order \(n\) the ratio is
\((\pi/2n)/\sin(\pi/2n)\):

\begin{center}
\begin{tabularx}{\textwidth}{@{}L{0.30\textwidth}Y Y Y Y@{}}
\toprule
Butterworth order \(n\) & 1 & 2 & 3 & 4 \\
\midrule
\(B_n/f_{\SI{3}{\decibel}}\) & 1.571 & 1.111 & 1.047 & 1.026 \\
Penalty (dB) & 1.96 & 0.46 & 0.20 & 0.11 \\
\bottomrule
\end{tabularx}
\end{center}

Sharper skirts make nameplate and noise bandwidth converge---which is a second,
independent reason to add poles, alongside the stopband slope of \cref{sec:filterspecs}.

Finally, the question of how wide to make the filter answers itself. Signal power
collected saturates once the passband covers the signal's own occupied bandwidth, while
noise keeps growing as \(B\):
\[
\mathrm{SNR}(B)=\frac{S(B)}{N_0B}
\qquad\Longrightarrow\qquad
\text{maximum near } B\approx B_{\text{signal}} .
\]
Too wide costs \(10\log_{10}(B/B_{\text{signal}})\) in pure noise; too narrow starts
truncating the signal. That is why receiver filters are specified per mode, and the
usual table of recommended bandwidths stops being a list to memorize.

\begin{exambox}
The chain of reasoning here is worth carrying whole, because most questions ask for one
link in it. The floor is \SI{-174}{dBm} per hertz, plus \(10\log B\) for bandwidth,
plus the noise figure in decibels---so \emph{narrower is quieter} and MDS follows
directly. In a cascade the first stage dominates and any loss ahead of the receiver
costs its full value in noise figure, which is the whole argument for a mast-mounted
preamplifier. At the top end, second-order products rise \SI{2}{\decibel} per
\si{\decibel} and third-order \SI{3}{\decibel} per \si{\decibel}, which makes
\(\mathrm{DR3}=\tfrac23(\mathrm{IIP3}-\mathrm{MDS})\) while blocking range is a plain
subtraction. An attenuator improves intermodulation by twice what it costs in signal.
And excessive local-oscillator phase noise raises the apparent floor near strong
signals---reciprocal mixing---which is a \(Q\) problem, not a distortion problem.
\end{exambox}
